\documentclass[10pt]{article}
% Shared presentation layer for the standalone R0--R7 development documents.
% Method equations remain authoritative in the canonical idea sketch.
\usepackage[a4paper,margin=18mm,headheight=14pt]{geometry}
\usepackage{fontspec}
\usepackage{xeCJK}
\xeCJKsetup{CJKspace=true}
\setmainfont{DejaVu Serif}
\setsansfont{DejaVu Sans}
\setmonofont{DejaVu Sans Mono}
\setCJKmainfont{Noto Serif CJK KR}
\setCJKsansfont{Noto Sans CJK KR}
\setCJKmonofont{Noto Sans Mono CJK KR}

\usepackage{amsmath,amssymb}
\usepackage{booktabs,longtable,tabularx,array}
\usepackage{enumitem}
\usepackage{xcolor}
\usepackage{hyperref}
\usepackage{fancyhdr}
\usepackage{microtype}
\usepackage[most]{tcolorbox}

\definecolor{CoreBlue}{HTML}{1F4E79}
\definecolor{CoreLight}{HTML}{EAF3F8}
\definecolor{StopRed}{HTML}{9C0006}
\definecolor{StopLight}{HTML}{FCE8E6}
\definecolor{GateGreen}{HTML}{38761D}
\definecolor{GateLight}{HTML}{EAF4E2}
\definecolor{DecisionOrange}{HTML}{B45F06}
\definecolor{DecisionLight}{HTML}{FFF2CC}

\hypersetup{
  colorlinks=true,
  linkcolor=CoreBlue,
  urlcolor=CoreBlue
}

\setlength{\parindent}{0pt}
\setlength{\parskip}{0.3em}
\setlength{\emergencystretch}{3em}
\setlist[itemize]{leftmargin=1.45em,itemsep=0.1em,topsep=0.2em}
\setlist[enumerate]{leftmargin=1.75em,itemsep=0.14em,topsep=0.22em}
\renewcommand{\arraystretch}{1.15}
\newcolumntype{L}[1]{>{\raggedright\arraybackslash}p{#1}}
\newcommand{\code}[1]{{\small\nolinkurl{#1}}}
\newcommand{\must}[1]{\textcolor{CoreBlue}{\textbf{#1}}}
\newcommand{\haltitem}[1]{\textcolor{StopRed}{\textbf{#1}}}
\newcommand{\passitem}[1]{\textcolor{GateGreen}{\textbf{#1}}}
\newcommand{\decisionitem}[1]{\textcolor{DecisionOrange}{\textbf{#1}}}

\newenvironment{contractbox}[1]
  {\begin{tcolorbox}[colback=CoreLight,colframe=CoreBlue,title={#1},fonttitle=\bfseries,before upper=\raggedright]}
  {\end{tcolorbox}}
\newenvironment{decisionbox}[1]
  {\begin{tcolorbox}[colback=DecisionLight,colframe=DecisionOrange,title={#1},fonttitle=\bfseries,before upper=\raggedright]}
  {\end{tcolorbox}}
\newenvironment{stopbox}[1]
  {\begin{tcolorbox}[colback=StopLight,colframe=StopRed,title={#1},fonttitle=\bfseries,before upper=\raggedright]}
  {\end{tcolorbox}}

\usepackage{fvextra}
\usepackage{xurl}
\DefineVerbatimEnvironment{verbatim}{Verbatim}{breaklines=true,fontsize=\scriptsize}
\hypersetup{pdftitle={Wind3DGS R1 Teacher 개발·실험 연구 기록},pdfauthor={Wind3DGS Project}}
\pagestyle{fancy}\fancyhf{}
\lhead{\small Wind3DGS R1 연구 기록}\rhead{\small 2026-09-09}\cfoot{\thepage}
\title{\textbf{Teacher 개발·실험 연구 기록}\\Registry부터 개발용 데이터와 공간 수렴 보완까지}
\author{Wind3DGS Project}\date{2026-09-06--09 작업; 2026-09-09 통합}
\begin{document}
\maketitle

\begin{contractbox}{이 기록의 용도와 판정}
이 채팅에서 이어 온 TeacherPhysicsRegistry, 저장·재생·probe, 사용자 GPU 검사, 구조 모델 진단,
시간·공간 수렴, 개발용 샘플 생성 및 P3 보완을 논문 작성 시 추적하기 위한 동반 기록이다.
최종 method authority는 상위 sketch, R1 acceptance 소유자는 \code{r1_teacher_probe_oracle.tex}다.
구현과 수치의 원본은 \code{../../code/}, \code{../../experiments/}에 있다.
\textbf{현재는 개발용 데이터 15개를 생성·검증한 상태이며, 본 학습 Teacher와 R1 전체는 미승인이다.}
P3의 처방 압력 공간 진단 통과와 원래 초기 변위 속도 진단 실패를 모두 유지한다.
\end{contractbox}

\tableofcontents\newpage

\section{연구 질문과 작업 범위}

Static 3DGS에서 wind-driven passive response package를 학습하려면 먼저 mesh teacher의 응답이
discretization이나 적분기 오차에 지배되지 않는지 확인해야 한다. 이 작업은 training-only reference를
만드는 과정이다. Target-mesh-free inference, student의 positive area/mass 및 Global/Local package
구조는 바꾸지 않았다. 고차 판 자체의 신규성을 이 작업의 논문 contribution으로 주장하지 않는다.

출발점은 \code{code/sessions/2026-09-06_01_teacher_dataset_handoff.md}의 기본 mesh/GUI였다.
후속 요청에 따라 registry, 정량 저장·비교, 실패 진단, 개발 sample까지 단계적으로 구현했다.
GUI 이후의 주된 확장은 물리 reference를 검증하고 데이터 provenance를 보존하는 CPU/CLI 경로다.
기존 Newton viewer가 새 P3 backend나 모든 정량 보고서를 자동으로 표시하는 것은 아니다.
그래프는 sample response, 공간 수렴 곡선과 내부 힘/모드 비교를 별도 artifact로 보존한다.

사용자는 필요한 선택에는 질문하되 샘플 생성·검증까지 연속 진행하도록 승인했다.
공간 보완 후 첫 accepted 데이터의 입력 범위(rest-start wind/on-off 우선인지, displaced free-decay까지
동시에 닫을지)는 질문했지만 확정 답변이 없다. 문서 정리 요청을 그 선택에 대한 승인으로 해석하지 않는다.

\section{구현 계보와 provenance 읽는 법}

Workspace 상대 경로 \code{code/}, \code{ideas/}, \code{experiments/}는 독립 저장소다.
GPU checkpoint는 code \code{7010682}, experiments \code{61d972e}에 이미 보존돼 있었다.
그 뒤 구조·수렴·샘플·보완 작업을 이번 checkpoint에서 함께 보존한다.
최종 commit 연결은 양쪽 \code{sessions/2026-09-09_01_teacher_checkpoint.md}와
ideas \code{sessions/2026-09-09_01_teacher_evidence_integration.md}를 따른다.
Run 생성 시에는 dirty source의 파일별 content hash가 실제 executable provenance다.
나중에 만든 commit ID가 당시 run의 source hash를 대신하지 않는다.

연구 기록의 수치는 각 report의 단위와 분모를 따른다. 초기 GPU fine-run RMS 상대 오차와
후속 고정 \(A,A\omega_*\) scale 오차는 다른 metric이며 직접 하나의 개선율로 합치지 않는다.
각 단계의 ``관련 테스트 수''는 범위가 겹치므로 누적 합산하지 않는다.
정적 fixture, 시간 오차, 공간 오차, replay, 데이터 무결성, R1 acceptance도 별도 판정이다.

\section{TeacherPhysicsRegistry와 데이터 경로}

\subsection{Registry의 기능과 제외 범위}

Pure schema와 Newton adapter를 분리했다. Native material parameter의 SI 해석, rest mesh와
source-object identity, material/solver/traction 설정, realized model 검사를 제공한다.
Registry 생성 자체가 simulation이나 파일 저장을 수행하지 않는다.
\code{validate_against_simulation()} 및 \code{require_valid()}로 실제 mass, pin,
material, rest geometry, gravity, solver 설정과 finite 상태를 검증한다.
Native bending의 단위 N을 연속 판 강성 N\,m로 바꿔 표기하지 않는다.

Strict JSON은 unknown/missing field, duplicate key, non-finite, unit/shape/version 오류를 거부한다.
Float 표현과 signed zero를 정규화하며 배열 hash는 little-endian C-order의 dtype/shape/unit/content를
포함한다. Effective registry hash, Teacher/GS 공용 traction identity hash, requested/run provenance와
실제 material array/solve 경로의 realized hash를 구분한다.
비활성 명목 설정과 시간 변화 wind speed를 모두 material identity로 섞지 않는다.
Source membership validator는 외부에서 검증한 split ref를 받으며 split 생성·배정·봉인을 대신하지 않는다.
Registry는 teacher-only이고 target package serializer를 새로 만든 것이 아니다.

초기 실제 model 대조는 \(\mathrm{rtol}=5\times10^{-6}\), mass/area absolute \(10^{-8}\),
pin \(10^{-6}\) m를 썼다. 이는 물리 수렴 tolerance가 아니다.
CPU/GPU solve branch와 환경을 기록하며 registry hash 동일성이 장치 간 bitwise trajectory 동일성을
보장하지 않는다. Core import는 Newton/Warp 없이 가능하도록 optional dependency 경계를 검사했다.

\subsection{Trajectory, sequence, 초기 상태}

Trajectory는 \(T\) interval과 \(T+1\) state를 구분하고 frame zero, timestamp, position/displacement,
velocity, held aerodynamic force와 interval work를 저장한다. Chunk inventory와 checksum으로 prefix를
검증하고 중간 실패의 완료 prefix, 오류와 미실행 부분을 보존한다. Reader의 무결성 검증과 실제 solver
replay는 별개다. 저장값만 다시 읽어 replay 통과로 표시하지 않는다.

Wind sequence runner는 초 단위 steady/pulse/chirp 등 프로그램을 case별 독립 초기 상태로 실행하고
성공·실패·미실행 inventory를 남긴다. Displaced free-decay 지원은 authored rest를 바꾸지 않고
초기 변위를 적용하며 requested 변위와 float32 realized state를 구분한다. Trajectory v2로 확장하고
v1 읽기 호환을 유지했다. Zero ambient는 상대 속도 drag가 남고, aero-off만 공력을 제거한다.
Material damping off 여부는 또 별도 설정이다.

\subsection{Common probe와 비교기}

Rest surface의 fixed probe, positive area/mass quadrature와 barycentric support를 저장한다.
Partition, constant/affine reproduction, coverage, unsupported reason 및 map hash를 검사했다.
현재 public map은 선형 3-support/nonnegative 계약이다. \(S=T\otimes I_3\)일 때
\[
 u_P=Su_T,\quad v_P=Sv_T,\qquad
 f_T=S^\top f_P,\quad f_T=S^\top W_A\tau_P,
 \qquad f_T^\top u_T=\tau_P^\top W_ASu_T.
\]
Total force N과 traction N/m\(^2\)를 구분하고 area를 두 번 곱하지 않는다.
Probe trajectory를 raw manifest에 연결하고 v1/v2 원본 재계산을 지원한다.
비교기는 같은 object/material/rest/input/BC 등 고정 조건을 검사한 뒤 velocity/tip/work/대역 PSD를
SI 및 무차원 값으로 보고한다. 실행 성공만으로 물리 수렴을 선언하지 않는다.
Independent GS 두 개의 reconstruction, common-valid intersection과 network-free oracle는 미완료다.

\section{사용자 GPU smoke: 실행 가능성과 수렴의 분리}

사용자가 \code{bash code/scripts/check_teacher_gpu.sh}를 실행했다.
GTX 1080 Ti 11 GiB, sm\_61, Warp 1.17.0, CUDA Toolkit 12.9/Driver 13.0에서
wind 5개, decay 2개, compare 3개, replay 2개의 \textbf{12/12 단계}가 통과했다.
Raw/probe 7쌍이며 pin drift와 guard activation은 0이다. Agent는 원본을 검토·재계산했고
추가 GPU simulation을 수행한 것으로 보고하지 않았다.
Raw: \code{experiments/artifacts/runs/teacher_gpu_check/20260907_042524_687651}.

공간 velocity 상대 RMS는 mesh 2$\to$4 10.07\%, 4$\to$8 24.90\%; temporal substeps
4$\to$8 23.99\%, 8$\to$16 11.41\%; aero-off decay mesh 4$\to$8 193.25\%였다.
분모는 fine run의 질량·시간 RMS 속도다. 짧은 0.2 s smoke이고
\code{convergence_status=not_assessed}다. 이는 발산률이나 변위 오차가 아니다.
같은 초기 probe 변위인데 native bending energy가 mesh별로 다른 관찰이 다음 감사를 이끌었다.

\section{Native hinge와 면적 보정 후보의 탈락 근거}

\subsection{고정 native bending 계수}

설치된 Newton 1.3.0의 에너지 의미를 독립 NumPy로 계산했다. Flat XY/XZ square에
\(w=A(X/W)^2\), \(W=H=1\) m, \(A=0.01\) m, \(k_e=10\) N와
\(n=4,8,16,32\)를 사용했다. Boundary edge는 bending 합에서 제외한다.
\[
 E_b=\frac12 k_e\sum_{e\text{ interior}}l_e^0\theta_e^2,
 \qquad K_E=\frac{2E_b}{A^2}.
\]
\(K_E\)는 에너지 등가 scalar이지 nonlinear force/A나 tangent HVP가 아니다.
Strip 기울기 \(s_j\)의 해석식
\[
 E_b=\frac12 k_eH\sum_j[\arctan(s_{j+1})-\arctan(s_j)]^2,
 \qquad K_{\rm small}=\frac{4k_eH}{W^2}\frac{n-1}{n^2}
\]
과 최대 상대 차이 \(1.792\times10^{-15}\)였다.
Energy는 \(0.000374911,0.000218695,0.000117157,0.000060531\) J,
\(K_E\)는 약 \(7.498,4.374,2.343,1.211\) N/m로 줄었다.
Finest/coarsest 비는 0.161454다. Native 계수를 고정하면 같은 연속 변형의 강성이 사라지는
mesh 의존성이 있으므로 그 계수 자체를 연속 판 강성으로 채택할 수 없다.
아주 작은 진폭의 float32 입력 문제가 있어 asymptotic fixture는 이진 정확 진폭 \(2^{-14}\)로
분리했다. 기본 \(A=0.01\) 결과나 tolerance를 사후 완화한 것이 아니다.

\subsection{Area-weighted hinge}

후보 \(k_e=B_hl_e^0/(A_L+A_R)\)에서 \(B_h\)의 단위는 N\,m다.
이는 native coefficient N과 구분되며 자동으로 isotropic plate \(D\)가 되지 않는다.
7개 변형, 두 대각선, 두 곡률, 네 mesh의 112개 사례를 계산했다.
작은 곡률의 \(\pm45^\circ\) 원통 bending에서 방향 에너지 비가 3으로 접근했고,
\(n=32\)에서 2.9375001853이었다. 한 방향에 맞춘 scalar calibration으로 해소되지 않는다.
\textbf{등방 Teacher mapping으로는 채택하지 않았다.}

\section{Quadratic patch KL 기준과 nonlinear shell}

\subsection{정적 선형 판}

\code{flat_plate_quadratic_patch_v1}은 central triangle의 edge-neighbor BFS patch를
1--4 ring에서 찾고 rank 6인 첫 design을 쓴다.
\([1,s,t,\tfrac12s^2,st,\tfrac12t^2]\)를 SVD로 fit하며 relative cutoff \(10^{-12}\),
condition limit \(10^8\)이다. Local length scale을 되돌린 curvature operator \(C\)와
engineering curvature \(\kappa=[\kappa_{xx},\kappa_{yy},2\kappa_{xy}]^\top\)를 사용한다.
\[
 D=\frac{Eh^3}{12(1-\nu^2)},\quad
 D_b=D\begin{bmatrix}1&\nu&0\\\nu&1&0\\0&0&(1-\nu)/2\end{bmatrix},\quad
 E_b=\frac12\sum_t A_t(C_tw)^\top D_b(C_tw),\quad
 f_b=-\sum_t C_t^\top A_tD_bC_tw.
\]
HVP는 이 선형 operator에서 정확하다. \(\nu=0,0.3\), 12방향 cylinder와 twist/dome/saddle,
quartic/sine, \(n=4\ldots32\)의 408개 정적 사례가 당시 기준을 통과했다.
Affine null 3개와 일반 변형 finest energy 오차 0.77\% 미만을 확인했다.
후속 공간 진단은 이 정적 통과가 동적 consistency를 보장하지 않음을 보여 주었다.

\subsection{3D 구조 연산자}

\code{flat_shell_stvk_projected_curvature_v1}은 StVK membrane의 Green strain과
current central normal에 투영한 bending을 합친 CPU 진단이다.
정확한 gradient/HVP에 normal derivative와 geometric-stress 항을 포함했다.
Finite area는 \(10^{-8}\)보다 커야 하고 proper rigid rotation을 단순 normal-flip으로 거부하지 않는다.
Rest Hessian의 PSD와 일반 current tangent의 가능한 indefiniteness를 구분한다.

\(\nu=0,0.3\), \(h=0.01,0.001\) m, 네 mesh의 1,200개 prescribed-geometry 사례에서
objectivity, rigid motion, gradient/HVP finite difference와 선형 극한은 통과했다.
그러나 8개 total-energy ladder는 \(n=16\to32\)에서 오차 감소에 실패했다.
Membrane/bending 오차 상쇄와 thin coarse mesh의 artificial membrane energy를 따로 남겼다.
대표 thin case의 ratio는 n4에서 약 185, n32에서 0.0447이었다.
이 현상만으로 뒤의 rest-normal checkerboard 결함을 설명하지 않는다.

\section{시간 적분: 정밀도 결함과 시간 해상도의 분리}

\subsection{초기 dynamics와 독립 reference}

Triangle-lumped total mass, position-only pin, Newmark/GMRES/line search를 갖는 CPU shell solver를
만들고 reaction, EOM residual, energy/work 및 모든 성공 state를 기록했다.
19 rollout, 2,105 step과 독립 재실행은 수치 계약을 통과했지만 nonlinear velocity의
160$\to$320 시간 대조는 11.394\%로 실패했다.
독립 DOP853 reference는 자체 대조를 통과했다.
추가 refinement에서 full N2560 velocity 11.7491\%, short N5120 5.50795\%였고,
short N10240은 5번째 step line search에서 실패했다. XY 고주파 시간 오차와
가까운 위치값을 빼서 가속도를 얻는 precision cancellation을 구분했다.

\subsection{Acceleration-form Newmark}

\(\beta=1/4,\gamma=1/2\)인 같은 Newmark 식을 acceleration unknown으로 다시 풀었다.
\[
 x_{n+1}=x_{\rm pred}+\beta\Delta t^2a_{n+1},\quad
 v_{n+1}=v_{\rm pred}+\gamma\Delta t a_{n+1},\quad
 R=Ma_{n+1}+\nabla E(x_{n+1})-f_{n+1}.
\]
Undamped tangent는 \(M+\beta\Delta t^2H\)다. 가까운 두 position을
\(\Delta t^2\)로 나누는 잔차 평가를 피했다. 물리 stiffness/damping, tolerance를 바꾸거나
고주파 필터·숨은 viscosity를 넣지 않았다. 기존 실패를 재현하고 새 3,457 step을 완료했다.
Short 차이는 2.20620\%로 개선됐지만 full N2560 11.7491\%는 남아 시간 해상도 문제임을 확인했다.

\begin{longtable}{L{0.20\linewidth}L{0.30\linewidth}L{0.40\linewidth}}
\toprule 구간 & refinement & velocity 차이와 판정\\\midrule
Short \(T_*/20\) & N10240/20480/40960 & 2.20620\% $\to$ 0.612308\% $\to$ 0.154044\%; 통과\\
Full \(T_*\) & N20480/40960/81920 & 6.059758\% $\to$ 2.974723\% $\to$ 0.868281\%; 통과\\
\bottomrule\end{longtable}

Full fixture의 \(T_*=1.0719083180935054\) s,
\(\omega_*=5.861681639298341\) rad/s는 n4 forward의 첫 positive normal mode 기준이다.
초기 amplitude는 0.001 m다. Short run은 full 기준 N의 분할 구간이며 새 계산 3,072 step/run이다.
Full은 143,360 step/run, 두 독립 run을 수행했다. 각 143,363 state, 560 chunk,
1,576,960 iteration/trial vector를 보존하고 state chain, force/reaction, EOM, Newmark update,
correction/Armijo/HVP와 비교 보고서를 재검산했다. Finest energy drift 기준은 \(10^{-3}\)이다.
각 run 약 45.7--45.8분, inventory 약 2.183 GB, 두 run 합 약 4.37 GB였다.
CPU/I/O 경합을 포함하므로 순수 solver 성능 benchmark가 아니다.
\textbf{이 통과는 해당 기존 shell과 초기 상태의 시간 진단이며 P3 teacher의 시간 통과가 아니다.}

\section{선형 공간 응답 실패를 확정한 실험}

\subsection{입력·경계·측정 고정}

1 m XY square, \(E=10^6\) Pa, \(\nu=0.3\), \(h=0.01\) m,
\(M_{\rm ref}=0.1\) kg, \(n=4,8,16\), forward/backward/checkerboard 세 diagonal이다.
\(x=0\)의 position만 고정하고 slope는 자유여서 normal rigid rotation null 하나가 남는다.
Clamped cantilever와 다른 경계다. \(u_z^0=10^{-3}x^2\), \(v^0=0\),
gravity/aerodynamics/damping 모두 off이며 동일 physical time \(T_*\), 10,241개 시각을 쓴다.

기존 shell의 rest-normal HVP에서 \(K\)를 조립하고 lumped \(M\)를 사용했다.
Mass-whitened generalized modal basis 전체를 써서 cos/sin exact response를 계산했다.
초기 상태를 mesh별 모드에 맞춰 다시 고르거나 high mode를 버리지 않았다.
Eigenvalue 원본을 보존하고 \(10^{-9}\max|\lambda|\) cutoff의 rigid null 하나를 별도 처리했다.
Time integration error를 제거한 상태에서도 공간 실패가 남는지 보는 진단이다.

Overlay16의 각 square를 네 triangle로 나누고 degree-2 3점 quadrature를 써 3,072개 probe를 만들었다.
이 규칙은 squared P1 difference에 정확하며 P3/spline 비교에는 충분하지 않다.
XY에서 기존 XZ mapper로 보낼 때 \((x,y,z)\mapsto(x,-z,y)\) proper rotation을 사용하고
dyadic rest position의 float32 표현을 확인했다. 평가 후 field를 원래 frame으로 되돌렸다.
Velocity는 \(A\omega_*\), displacement는 \(A=0.001\) m로 normalize하고 finest 1\%와
오차 감소를 사전 기준으로 썼다. Numerical floor는 별도 \(10^{-10}\)이다.

\begin{longtable}{L{0.27\linewidth}L{0.28\linewidth}L{0.35\linewidth}}
\toprule 비교 & displacement [\%] & velocity [\%]\\\midrule
Forward/backward 4$\to$8 & 5.956244 & 22.169172\\
Forward/backward 8$\to$16 & 2.927207 & 23.209961\\
Checkerboard 4$\to$8 & 3.857861 & 21.043244\\
Checkerboard 8$\to$16 & 1.395752 & 22.685509\\
n16 forward/backward & 0.570938 & 8.558328\\
n16 forward(or backward)/checker & 21.203509 & 51.546194\\
\bottomrule\end{longtable}

Maximum energy drift \(6.0021\times10^{-12}\), EOM \(1.4526\times10^{-12}\)로
수치 계산은 통과했지만 공간/방향 기준은 실패했다. 각 run 92,169 frame, 720 chunk,
1,509 inventory 약 266 MB다. 독립 두 run에서 runtime 외 1,508개 파일이 동일했다.
이 실패 뒤 dataset을 확대해 물리 결함을 덮지 않고 개발 sample만 분리 생성했다.

\section{개발용 학습 sample 15개의 실제 생성과 검증}

기존 Newton native backend에서 CPU로 새 source 세 개를 만들었다.
1 m XZ square n8, \(M=0.1\) kg, fps60, substeps16, iterations10, 1 s/61 state,
5$\times$5 trapezoid probe 25개다. Native stretch/area stiffness \(10^3\), material damping 0.1,
bending 10 N, bending damping 0, gravity 0을 사용했다.
\begin{itemize}
 \item Wind pulse: 0.5 m/s, 0.2 s half-sine 뒤 0.8 s zero ambient.
 \item Step on/off: 0.5 m/s의 켜짐/꺼짐. Off에도 air drag가 남는다.
 \item Aero-off decay: \(\Delta Y=0.01x^2\), 공력 off, material damping 유지.
\end{itemize}
각 source의 60 interval을 12 interval/13 state window 5개로 나눠 총 15개를 만들었다.
Target array는 \([13,25,3]\), batch 크기는 \([4,4,4,3]\)이다.
Single development source-object group이며 train/dev/test 독립 객체 분할이나 GS pairing을 완료한 것이 아니다.

Raw/probe inventory와 label/input/work를 모두 대조하고 각 CPU source를 독립 재실행했다.
세 replay의 position/velocity/aerodynamic force/work 최대 차이는 모두 0,
pin drift와 guard count도 0이었다. 각각 maximum nodal speed는
0.202946126/0.471042246/0.133365527 m/s,
interval work 절댓값 합은 0.000789484374/0.006596602165/0 J다.
마지막 값은 순 외력 일이나 구조 energy drift가 아니다.
생성·대조·replay·batch 검증은 48.370 s가 걸렸다.

Dataset에는 static probe/weight/attachment, registry, split, per-window NPZ와 checksum manifest를 둔다.
\code{TeacherSampleDataset.open(..., allow_development=True)}로만 개발 데이터를 열 수 있고,
default는 \code{training_eligible=false} 데이터를 거부한다.
\code{verify_sources}는 manifest뿐 아니라 raw/probe에서 label을 재구성해 대조한다.
Tamper, incomplete output, rehashed corruption과 static geometry 오류를 검사했다.
저장된 NumPy batch 읽기는 Torch/Newton/Warp 없이 가능하다.

Raw source는 \code{experiments/artifacts/runs/teacher_sample_dataset/20260908_sample_v1},
dataset은 \code{experiments/artifacts/datasets/teacher_samples/20260908_sample_v1}다.
Source inventory 93개/1,150,138 byte, dataset manifest 포함 20개/125,293 byte다.
Compact evidence와 response PNG를 Git에 보존하고 NPZ 원본은 ignored artifact에 둔다.
\textbf{이 데이터는 P3 데이터가 아니며 후속 P3 통과로 eligibility를 올리지 않는다.}

\section{공간 수렴 보완: 원인과 대안 분리}

\subsection{Quadratic patch의 내부 weak-force 결함}

상수 quadratic curvature를 넣어도 checkerboard mesh의 안전한 내부 vertex에 약
\(9.5238095238\times10^{-4}\) N의 force가 n8/16/32에서 남았다.
경계에 닿는 patch의 영향을 받는 모든 vertex를 제외해 safe interior 9/121/729개를 검사했다.
Forward/backward는 약 \(10^{-17}\)--\(10^{-15}\) N의 roundoff 수준이었다.
Polynomial curvature/energy 재현과 그 adjoint가 만드는 약형 내부 힘의 consistency는 다른 조건이다.

Safe interior에만 \(O(h_{\rm mesh}^2)\) parity perturbation을 가해 에너지를 최소화하면
8.1794\%, 21.6639\%, 31.0786\% 감소했다. Amplitude는
\(3.495\times10^{-6},6.886\times10^{-7},1.640\times10^{-7}\) m였다.
작아지는 요철로도 mesh-dependent energy를 낮출 수 있는 반례다.
Ring을 2--4로 늘려도 잔류 force가 각각 약
\(2.97\times10^{-5},8.64\times10^{-5},6.11\times10^{-5}\) N로 남았다.
Patch ring 확대나 mass 교체만으로 해결됐다고 판단하지 않았다.

Lumped mass를 consistent mass로 바꾼 checker n16의 두 번째 positive frequency는
11.13933에서 11.239343 rad/s로 바뀌지만 독립 기준 14.259993과의 차이가 남는다.
Old n16의 첫 5개 positive frequency는 다음과 같다 (rad/s).
\[
\begin{aligned}
\text{forward}:&\quad[6.33265,14.18110,24.15372,24.60083,45.84505],\\
\text{checker}:&\quad[6.20080,11.13933,20.19455,21.38170,36.04130].
\end{aligned}
\]
Stiffness consistency와 mass approximation을 구분한 근거다.

\subsection{독립 spline와 C0 interior-penalty P2/P3}

독립 기준은 unit square의 tensor-product \(C^2\) cubic B-spline, 4/8/16/32 spans,
Gauss4 tensor quadrature의 \(H^2\)-conforming KL 판이다. Consistent \(M,K\)를 조립하고
x0 첫 coefficient만 제거해 slope-free 경계를 유지한다. Quadratic 초기 상태는 Greville collocation으로
정확히 표현하고 generalized eigenbasis 전체를 사용한다.
Simply-supported analytic frequency를 별도 구현 검증으로 사용했으며 주 비교의 경계를 바꾸지 않았다.

P2는 triangle당 6 DOF, P3는 10 DOF의 Lagrange shape다. P3는 local Vandermonde로 shape/gradient/Hessian을
계산한다. Volume KL energy에 interior edge moment consistency와 slope-jump penalty를 더한다.
\[
 E=\frac12\sum_t\int_t\kappa^\top D_b\kappa\,dA
 -\sum_e\int_e[\partial_n w]\langle M_{nn}\rangle\,ds
 +\frac12\sum_e\int_e\frac{\alpha}{h_{\rm avg}}[\partial_n w]^2\,ds.
\]
Normal orientation과 jump/average는 source에 고정돼 있다.
\(h_{\rm avg}\)는 이웃 cell diameter의 평균이며 물리 두께와 다른 길이다.
P2는 \(\alpha=Eh^3\), P3는 \(\alpha=Eh^3(p/2)^2\), \(p=3\)의 고정 factor 2.25다.
P3 결과를 본 뒤 penalty를 sweep하거나 유리한 값을 골라낸 것이 아니다.
P2의 volume/edge는 Duffy3/Gauss2, P3는 Duffy4/Gauss3를 사용한다.
Exterior rotation penalty는 추가하지 않았다.

식의 reference는 선행 구현 기록에 연결한 공식 FEniCS-Shells Kirchhoff--Love demo다.
해당 demo의 clamped 경계를 이번 position-only 경계로 복사하지 않았다.
Reference URL은 \code{experiments/R1_teacher_spatial_remediation/README.md}, 당시 적용 판단은
\code{code/sessions/2026-09-08_07_teacher_spatial_remediation.md}에 보존돼 있다.
이번 문서 통합에서는 외부 문헌을 새로 조회하지 않았으며 신규성 조사를 수행한 것으로 표시하지 않는다.

P2/P3는 constant-curvature interior force가 각각 약 \(10^{-15}/10^{-14}\) N,
unconstrained affine null 3개/left-pin null 1개, proper rotation, gradient, area/mass,
patch polynomial force/energy를 검사했다.
Independent oscillator matrix exponential에는 zero mode와 resonance도 포함했다.
P3 n16 forward frequencies는
\([6.35772073,14.26003426,24.28380007,24.88148287,46.36581911]\),
spline32는 \([6.35770535,14.25999311,24.28325170,24.88117461,46.36373132]\) rad/s다.

\section{두 입력군, 정확한 공간 norm과 연속 시간 상계}

원래 \(w^0=10^{-3}x^2,v^0=0\)를 유지하고, 별도 진단으로 rest-start uniform pressure
\(p(t)=0.01\sin(\pi t/0.2)\) Pa를 \(0\le t\le0.2\) s에 가했다.
그 뒤 pressure는 0이다. 이 prescribed pressure는 relative wind quadratic traction도,
canonical frame-start held aerodynamic force도 아니다. 새 wind dataset로 소개하지 않는다.

모든 finite modal equation \(q''+\omega^2q=g\sin(at)\), \(a=\pi/0.2\)를 정확히 계산하고,
zero/resonant mode에는 안정적인 sinc/극한식, off 구간에는 homogeneous solution을 썼다.
기존 10,241개 시각과 \(A=0.001\) m, \(A\omega_*\) scale을 그대로 사용했다.
P3--P3 difference는 overlay16 Duffy4로 degree-6 squared field를 적분한다.
P3--tensor cubic spline은 overlay32 Duffy6로 total degree 최대 9의 곱을 적분한다.
P1 전용 3,072 probe 규칙을 그대로 재사용하지 않았다.

Mass-normalized continuous surface RMS를 \(\|\cdot\|\)라 하면 sampled maximum만 보고
시각 사이 peak를 놓치지 않도록 다음 finite-modal continuous-time 상계를 추가했다.
\[
 \max_{0\le t\le T_*}\|e(t)\|
 \le \max_j\|e(t_j)\|+L_e\frac{\Delta t}{2},\qquad
 L_u\ge\max_t\|\dot e_u(t)\|,\quad L_v\ge\max_t\|\dot e_v(t)\|.
\]
Pair의 각 model velocity/acceleration bound를 합하고 consistent-mass orthonormal basis와
\(\sqrt{M_{\rm ref}}\) normalization을 썼다. Coarse sampled lower와 fine continuous upper로
refinement 감소도 검사한다. 이는 주어진 finite modal 모델의 시간 sample 사이 상계이며
미해결 continuum error나 floating-point uncertainty 전체의 수학적 인증은 아니다.

Pressure mode의 \(|v|\)에는 \(2|g|/a\), on-interval \(|q|\)에는
\(\min(|g|\tau^2/2, |g|(1+a/\omega)/|\omega^2-a^2|)\) 형태의 유효 bound를 쓰고,
zero/resonance에서는 특수 fallback을 사용한다. Acceleration은 on에서
\(|g|+\omega^2|q|\), off에서 \(\operatorname{hypot}(\omega^2q_\tau,\omega v_\tau)\)다.
고주파를 제거하거나 입력별로 norm 분모를 재조정하지 않았다.

\begin{longtable}{L{0.39\linewidth}L{0.20\linewidth}L{0.29\linewidth}}
\toprule Pressure 비교 & displacement [\%] & velocity [\%]\\\midrule
Spline 8$\to$16 sampled & 0.0176344 & 0.191323\\
Spline 16$\to$32 sampled & 0.00105483 & 0.0234081\\
P2 8$\to$16 sampled 최대 & 2.245875 & 6.079196 (실패)\\
P3 8$\to$16 sampled 최대 & 0.016764 & 0.136376\\
P3 8$\to$16 continuous upper 최대 & 0.149654 & 0.526612 (통과)\\
P3 n16 direction continuous upper 최대 & 0.133208 & 0.398941 (통과)\\
P3 n16/spline32 continuous upper 최대 & 0.134202 & 0.408232 (통과)\\
\bottomrule\end{longtable}

P2 direction velocity도 1.106636\%로 실패했다. P3의 pressure 공간/방향/기준 비교는 고정 1\%를
통과했다. P3 free-energy drift 최대는 \(1.0684320\times10^{-8}\), pressure의 trapezoid
work/energy maximum defect는 \(1.5848926\times10^{-7}\)이다.
Work 수치는 정확 연속 시간 work 적분을 계산한 값으로 표기하지 않는다.

반면 원래 \(x^2\) 초기 상태의 P3 8$\to$16 velocity sampled error는
17.5895--17.9281\%, continuous upper는 약 24.27--24.80\%로 실패한다.
Displacement sampled error는 0.3017--0.3042\%다.
Spline의 같은 입력 velocity는 4$\to$8 27.3165\%, 8$\to$16 19.1290\%,
16$\to$32 14.9829\%로 감소하지만 아직 1\%에 못 미친다.
Spline32 초기 에너지의 16.8013\%가 100 rad/s 위, 8.4490\%가 500 rad/s 위에 있다.
이 초기 상태는 유효한 finite-energy 상태이며 ``불법 초기 조건''으로 제외하지 않는다.
정상 pressure 응답과 broad-band initial velocity의 수렴 난이도 차이로 해석하고 추가 검증 대상으로 둔다.

\clearpage
\section{스케치에 직접 반영한 계약 보정과 보류}

\begin{longtable}{L{0.23\linewidth}L{0.33\linewidth}L{0.34\linewidth}}
\toprule 항목 & 반영한 판단 & 아직 완료하지 않은 것\\\midrule
Teacher sample/DOF & 고차 DOF와 positive quadrature를 구분. \(M=(M_{\rm ref}/A_{\rm ref})\Psi^\top W_A\Psi\) 허용 & 최종 registry의 고차 law 및 actual model validator\\
Map sign & GS/기존 P1은 nonnegative 유지. 고차 Teacher에만 signed shape와 별도 version 허용 & P3 10-shape public map, coverage/noise/adjoint와 trajectory 연결\\
Convergence identity & 같은 backend/mass/BC/load에서 spatial/time acceptance 요구 & P3 nonlinear wind의 temporal+spatial+FRF/tip/work 검증\\
학습 적격성 & Development sample을 명시적 opt-in으로 격리 & Accepted Teacher와 source split, independent GS 및 oracle\\
첫 input family & Rest-start pressure 진단은 별도 scope & Wind/on-off 우선 여부의 사용자 선택; 기존 decay 요구 유지\\
Claim & P3는 training-only reference 후보, target method 변경 없음 & Learned held-out 개선이나 full Teacher 적격성 주장 불가\\
\bottomrule\end{longtable}

Full shell 연결에는 objective membrane+bending, finite rotation, consistent mass,
attachment/reaction, nonlinear tangent/solve, damping, canonical two-sided relative wind,
guard-before-area, frame-start hold와 work quadrature가 필요하다.
선형 P3 pressure 모듈을 기존 3-node/lumped-mass nonlinear solver에 이름만 바꿔 넣지 않는다.
같은 새 backend에서 해상도/시간 검증을 다시 수행한 뒤 producer를 연결해야 한다.
Independent GS 두 개, common mask, spectrum convergence와 network-free oracle/base transport는 그 후에도
별도 R1 종료 조건이다. Student 학습이나 대규모 dataset 생성은 아직 수행하지 않았다.

\section{재현 command, hash와 원본 보존}

명령은 workspace root에서 실행한다. Launcher가 \code{code/}로 이동하므로 launcher 인자 경로는
code 기준이다. 모든 새 run은 사용하지 않은 output 이름을 요구한다. 기존 artifact를 덮어쓰거나
자동 재개하지 않는다. 전체 물리 재실행과 inventory/source/gate 검산을 구분한다.

\begin{verbatim}
bash code/scripts/check_teacher_gpu.sh
bash code/scripts/generate_teacher_sample_dataset.sh \
  --output ../experiments/artifacts/runs/teacher_sample_dataset/my_run \
  --dataset ../experiments/artifacts/datasets/teacher_samples/my_dataset
bash code/scripts/audit_teacher_plate_spatial_remediation.sh \
  --source-run ../experiments/artifacts/runs/teacher_shell_linear_spatial/20260908_reference_v1 \
  --output ../experiments/artifacts/runs/teacher_spatial_remediation/my_reference
bash code/scripts/audit_teacher_plate_cubic_refinement.sh \
  --source-run ../experiments/artifacts/runs/teacher_spatial_remediation/20260908_reference_v1 \
  --output ../experiments/artifacts/runs/teacher_spatial_remediation/my_cubic
bash code/scripts/audit_teacher_plate_cubic_refinement.sh \
  --verify ../experiments/artifacts/runs/teacher_spatial_remediation/my_cubic
\end{verbatim}

전체 시간 refinement는 네 선행 run이 필요하다. 정확한 인자 및 재검산 명령은
\code{experiments/R1_teacher_shell_full_refinement/README.md}에 있다.
다른 단계의 launcher/인자도 아래 evidence index의 README가 소유한다.
P3의 \code{--verify}는 inventory/current source/판정 검산이며 전체 modal 물리 재계산은
별도 fresh run으로 실행했다.

\begin{longtable}{L{0.27\linewidth}L{0.65\linewidth}}
\toprule Artifact & SHA-256 (전체 값)\\\midrule
Full temporal report & \code{14700c5346b3f9502f724abd32febcc7a7a6266771ab08fbb750dae531e985f5}\\
Linear spatial report & \code{3a0d8e7ee5be505b190008d387c67a113a027ce0651595cb43b0812c8c33ba1a}\\
Dataset manifest & \code{1ab229207fa8ec60c282aa433c49d7a5f8a288deb1575688794b72b2afdae8f5}\\
Remediation reference report & \code{b5c3929f2e65160f9f8127fe4ea9612e41a0871a5554c4efd5bcec0e77249fd7}\\
P3 report & \code{799a48aa00e9eb40085ad030298d90a409959332092fbfdc784b312bc2cb46ac}\\
\bottomrule\end{longtable}

Remediation raw root는 \code{experiments/artifacts/runs/teacher_spatial_remediation/}다.
Reference는 \code{20260908_reference_v1}와 \code{20260908_reference_v1_replay},
P3는 \code{20260908_cubic_v1}와 \code{20260908_cubic_v1_replay}다.
Reference run은 각각 45.784/41.826 s, inventory 55개/약 45.16 MB;
P3 run은 104.291/106.025 s, inventory 16개/146,317,681 byte다.
각 pair는 runtime 외 각각 54/15개 inventory가 byte-identical이다.
모델 배열과 전체 comparison curve는 독립 계산으로 재생성했다.
모든 native nodal trajectory를 따로 저장한 것은 아니며 보존한 modal model과 계수로 재구성한다.
Preceding spatial source 30개와 새 source snapshot, config/environment 및 모든 inventory hash를 연결한다.

Git에는 compact report/config/environment/provenance/verification과 선택한 PNG를 보존한다.
Ignored raw의 모든 state/trace/vector, modal arrays 및 dataset NPZ는 Git clone만으로 복구되지 않는다.
Raw artifact를 별도 보존하거나 기록된 source/config로 재생성해야 한다. Compact evidence만으로
NPZ가 복원된다고 주장하지 않는다. 개인 절대 경로·secret은 기록에서 제외한다.

\section{전체 상세 기록 index와 논문 작성 시 사용법}

Code 기능별 설계·승인·수정·검증 과정은 아래 session에 남아 있다. 날짜의 각 번호는 해당 repository
내 순서다. 2026-09-01 viewer 기록은 인수한 선행 맥락이며 이후 작업과 구분한다.
\begin{longtable}{L{0.20\linewidth}L{0.70\linewidth}}
\toprule Code session 날짜/번호 & 소유 내용\\\midrule
09-06 01, 02 & Dataset handoff; TeacherPhysicsRegistry native SI/schema/실제 model 검사\\
09-07 01, 02, 03 & Trajectory writer/reader/replay; wind sequence; authored-rest 초기 변위\\
09-07 04, 05, 06, 07 & Common probe/adjoint; convergence comparator; GPU launcher; GPU checkpoint\\
09-07 08, 09, 10 & Native bending 감사; area-weighted hinge; quadratic patch plate 대안\\
09-07 11, 12 & 3D shell structure; CPU nonlinear dynamics\\
09-08 01, 02, 03, 04 & Temporal reference/실패; acceleration precision; short 및 full refinement\\
09-08 05, 06, 07 & Linear spatial 실패; sample 생성·검증; stencil/P2/P3 보완\\
09-09 01 & 누적 코드 회귀·Git checkpoint와 현재 다음 작업\\
\bottomrule\end{longtable}

\code{experiments/} 아래 다음 각 폴더의 \code{README.md}에는 해당 입력, 결과 표, raw 경로,
명령, source/config hash와 재검산 근거가 있다. 별도의 experiment session은 실행 과정과 실패를 보존한다.
\begin{longtable}{L{0.47\linewidth}L{0.43\linewidth}}
\toprule Evidence folder & 논문에서 참고할 근거\\\midrule
\code{R1_teacher_smoke} & 사용자 GPU 실행과 물리 수렴 미판정\\
\code{R1_teacher_bending_audit} & Native coefficient의 mesh scaling 반례\\
\code{R1_teacher_bending_mapping} & Area hinge의 방향 비등방성\\
\code{R1_teacher_plate_reference} & Static KL polynomial/방향 검증의 범위\\
\code{R1_teacher_shell_structure} & Objectivity/미분, 오차 상쇄와 thin locking 진단\\
\code{R1_teacher_shell_dynamics} & 수치 계약 통과와 초기 시간 실패\\
\code{R1_teacher_shell_temporal} & 독립 DOP853, 고주파와 line-search 실패\\
\code{R1_teacher_shell_precision} & Acceleration unknown의 정밀도 보정 근거\\
\code{R1_teacher_shell_refinement} & Short 시간 해상도 ladder\\
\code{R1_teacher_shell_full_refinement} & Full 시간 수렴·전체 trace 검산\\
\code{R1_teacher_shell_linear_spatial} & Exact temporal response에서도 남는 공간/방향 실패\\
\code{R1_teacher_sample_dataset} & 실제 15 window, source/replay/loader 검증\\
\code{R1_teacher_spatial_remediation} & 내부 force 반례, 독립 spline, P2/P3와 입력별 차이\\
\bottomrule\end{longtable}

논문 method에는 mass/shape/quadrature, attachment와 force sampling을 정확히 옮기고,
appendix에는 실패한 native/area-hinge/patch와 보완 근거를 선택해 인용할 수 있다.
Results에서는 실행 성공, numerical convergence와 accepted Teacher를 분리하고 normalized error의
분모와 측정 시각/연속 상계 여부를 함께 쓴다. Limitations에는 broad-band displaced response,
미완성 nonlinear wind 연결, 독립 GS/oracle 및 R1 전체 미완료를 남긴다.
이 기록 자체는 아직 수행하지 않은 learner 실험이나 physics validation의 대체 증거가 아니다.

\end{document}
